\chapter{Functions}

\section{Why This Matters}
A function is a reusable block of code that performs a specific task. You "define" it once, and then you can "call" it as many times as you want. This prevents copying and pasting code, adhering to the \textbf{DRY} (Don't Repeat Yourself) principle.

In Machine Learning, functions are used to:
\begin{itemize}
    \item Load data (\texttt{load\_dataset()})
    \item Preprocess data (\texttt{normalize\_image(image)})
    \item Train models (\texttt{train\_step(model, data, labels)})
\end{itemize}

\section{Explanation}
You define a function using the \texttt{def} keyword, followed by the function name, parentheses \texttt{()}, and a colon \texttt{:}. 

\subsection{1. Inputs (Parameters vs. Arguments)}
The variables listed in the definition are called \textbf{parameters}. The actual data you pass in when calling the function are called \textbf{arguments}.

\begin{lstlisting}
def greet(name):          # 'name' is the parameter
    print(f"Hello, {name}!")

greet("Alice")            # 'Alice' is the argument
\end{lstlisting}

\subsection{2. Outputs (The \texttt{return} Statement)}
If you want a function to calculate a value and \textit{hand it back} to you, you must use the \texttt{return} keyword. 

\warning{Once a function hits a \texttt{return} statement, it exits immediately. No code below the \texttt{return} will run!}

\begin{lstlisting}
def multiply(a, b):
    result = a * b
    return result
    print("This will never print!") # Unreachable code

answer = multiply(5, 10)
\end{lstlisting}

\subsection{3. Returning Multiple Values}
In Python, you can return multiple values separated by commas. Under the hood, Python secretly packages them into a \textbf{Tuple}! You can use Tuple Unpacking to catch them.

\begin{lstlisting}
def get_model_stats():
    accuracy = 0.95
    loss = 0.05
    return accuracy, loss  # Secretly a tuple!

acc, lss = get_model_stats()
\end{lstlisting}

\mlconnection{When you write a custom neural network layer or an evaluation script, you will always wrap it in a function. For example, a standard evaluation function in PyTorch returns multiple values like \texttt{final\_accuracy} and \texttt{final\_loss}.}

\section{Solutions to Exercises}

\textbf{Exercise 1: The Basic Calculator}
\begin{lstlisting}
def subtract(x, y):
    return x - y
print(subtract(10, 3))
\end{lstlisting}

\vspace{1em}
\textbf{Exercise 2: The Data Normalizer}
\begin{lstlisting}
def normalize_pixel(pixel_value):
    return pixel_value / 255.0
print(normalize_pixel(128))
\end{lstlisting}

\vspace{1em}
\textbf{Exercise 3: Early Return Logic}
\begin{lstlisting}
def check_access(role):
    if role == "admin":
        return True
    return False
\end{lstlisting}
\textit{Solution: Notice the lack of an \texttt{else} block! Because the \texttt{return} statement instantly exits the function, if the role is "admin", it returns True and dies. If it is NOT "admin", the if block is skipped, and it hits the \texttt{return False} at the bottom. This is called "Early Return".}
